% Source: https://artofproblemsolving.com/community/c3h1731721_2006_mathcounts_national_sprint_round_21
% Stable identity: aops-topic-1731721-q01
``Indistinguishable'' means that the two T's are identical, so collections that differ only in which T was taken are the same collection; we count collections by which letters they contain, not by which magnets.

The vowels are A, O, U (all different): $\binom32=3$ choices of two.

The consonants are M, T, H, C, N, T, S: six different letters, with T appearing twice. Choose three by the number of T's:
\begin{itemize}
\item no T: $\binom53=10$ (from M, H, C, N, S);
\item one T: $\binom52=10$;
\item two T's: $\binom51=5$.
\end{itemize}
That gives $25$ consonant collections. In total there are $3\cdot25=\boxed{75}$ collections.
